\documentclass[12pt]{article}

\usepackage{amsmath, amssymb, bm}
\usepackage{booktabs}
\usepackage{geometry}
\usepackage{setspace}
\geometry{margin=1in}
\doublespacing

\title{Variance Components from Nuclear Family Trios:
A Linear Mixed Model Implementation in R}
\author{}
\date{}

\begin{document}
\maketitle

\section{Introduction}

Quantitative genetic variance can be decomposed into additive genetic (A),
dominance genetic (D), shared environmental (C), and unique environmental (E)
components.

Twin studies traditionally estimate ACE or ADE models, but many modern
cohorts consist of nuclear family trios (mother--father--child).
This paper presents a linear mixed model framework for estimating AE, ACE,
and ADE models from trio data using nlme in R.

\section{Variance Decomposition}

For an individual phenotype $y$,

\[
y = \mu + a + d + c + e
\]

\[
Var(y) = \sigma_A^2 + \sigma_D^2 + \sigma_C^2 + \sigma_E^2
\]

\section{General Pedigree Covariance}

The covariance between relatives $i$ and $k$ under the ACDE model is

\[
\text{Cov}(y_i,y_k)
=
2\phi_{ik}\sigma_A^2
+ \Delta_{ik}\sigma_D^2
+ \gamma_{ik}\sigma_C^2
\]

\begin{itemize}
\item $\phi_{ik}$ = kinship coefficient (additive sharing)
\item $\Delta_{ik}$ = probability of sharing two alleles IBD (dominance)
\item $\gamma_{ik}$ = indicator of shared family environment
\end{itemize}

\section{Coefficients for Nuclear Family Trios}

\begin{center}
\begin{tabular}{lccc}
\toprule
Pair & $2\phi$ & $\Delta$ & $\gamma$ \\
\midrule
Mother--Father & 0 & 0 & 1 \\
Mother--Child  & $\tfrac12$ & 0 & 1 \\
Father--Child  & $\tfrac12$ & 0 & 1 \\
\bottomrule
\end{tabular}
\end{center}

\subsection{Key Consequence}

Dominance deviations are not transmitted from parents to offspring:

\[
\Delta_{ik}=0 \quad \text{for all trio pairs}
\]

Therefore dominance contributes only to individual variance:

\[
Var(y)=\sigma_A^2+\sigma_D^2+\sigma_E^2
\]

and produces no observable covariance between relatives.

\section{Identifiable Family Resemblance}

Trio data contain only one source of familial covariance:

\[
\text{Cov}_{family} =
\frac12\sigma_A^2 + \sigma_F^2
\]

where $\sigma_F^2$ is an unidentified family resemblance component.

This component can be interpreted as:

\[
\sigma_F^2 =
\begin{cases}
\sigma_C^2 & \text{ACE interpretation} \\
\text{residual familial resemblance} & \text{ADE interpretation}
\end{cases}
\]

Hence ACE and ADE models are mutually exclusive in trio data.

\section{Mixed Model Representation}

We model stacked trio observations:

\[
\mathbf{y} = \mathbf{X}\beta + \mathbf{Z}_A \mathbf{a} + 
\mathbf{Z}_F \mathbf{f} + \boldsymbol{\varepsilon}
\]

\begin{align}
\mathbf{a} &\sim N(0, \sigma_A^2 \mathbf{A}) \\
\mathbf{f} &\sim N(0, \sigma_F^2 \mathbf{I}) \\
\boldsymbol{\varepsilon} &\sim N(0, \sigma_E^2 \mathbf{I})
\end{align}

\section{Data Restructuring in R}

\begin{verbatim}
library(nlme)

longdata <- reshape(
  widelong,
  varying = c("bmi1","bmi2","bmi3"),
  v.names = "bmi",
  timevar = "member",
  times = c("mother","father","child"),
  idvar = "famid",
  direction = "long"
)

longdata$Acoef <- ifelse(longdata$member=="child",1,0.5)
longdata$Dcoef <- 1
\end{verbatim}

\section{AE Model}

\[
Var(y) = \sigma_A^2 + \sigma_E^2
\]

\begin{verbatim}
model_AE <- lme(
  bmi ~ 1,
  random = list(
    famid = pdDiag(~ Acoef - 1),
    indid = pdDiag(~ Dcoef - 1)
  ),
  data = longdata,
  method = "REML"
)
\end{verbatim}

\section{ACE Model}

\[
Var(y) = \sigma_A^2 + \sigma_C^2 + \sigma_E^2
\]

\begin{verbatim}
model_ACE <- lme(
  bmi ~ 1,
  random = list(
    famid = pdBlocked(list(
      pdDiag(~ Acoef - 1),
      pdDiag(~ Dcoef - 1)
    )),
    indid = pdDiag(~ Dcoef - 1)
  ),
  data = longdata,
  method = "REML"
)
\end{verbatim}

\section{ADE Model}

\[
Var(y) = \sigma_A^2 + \sigma_D^2 + \sigma_E^2
\]

In trio data the ADE model estimates additive genetic variance and a
combined non-additive/unique environmental component.

\begin{verbatim}
model_ADE <- lme(
  bmi ~ 1,
  random = list(
    famid = pdDiag(~ Acoef - 1),
    indid = pdDiag(~ Dcoef - 1)
  ),
  data = longdata,
  method = "REML"
)
\end{verbatim}

\section{Extracting Variance Components}

\begin{verbatim}
VarCorr(model_ACE)
VarCorr(model_AE)
VarCorr(model_ADE)
\end{verbatim}

\section{Model Comparison}

\begin{verbatim}
anova(model_AE, model_ACE)
anova(model_AE, model_ADE)
\end{verbatim}

\section{Interpretation}

\begin{itemize}
\item AE vs ACE tests shared environmental effects.
\item AE vs ADE tests non-additive genetic effects.
\item ACE vs ADE cannot be fit simultaneously.
\end{itemize}

\section{Conclusion}

Nuclear family trios allow estimation of additive genetic variance and a
single familial resemblance component. This component may be interpreted
as shared environment (ACE) or as residual familial resemblance (ADE).
Model comparison provides evidence for the most plausible interpretation.

\end{document}
